\section{Measure the isolated kernel before changing the codec}
The verified construction sources produce the following entry-to-return counts for
4096 native words. They retain the observation labels and their \texttt{nop}
instructions. Input preparation, assembly, console printing, record encoding, CRCs
and display timing are outside this interval.
\Needspace{7\baselineskip}
\begin{longtable}{p{30mm}p{30mm}p{38mm}p{43mm}}
\toprule Mode & Native words & Simulated instructions & Instructions per word \\\midrule
Scalar & 4096 & 53\,269 & 13.005 \\
MSA & 4096 & 29\,721 & 7.256 \\\bottomrule
\end{longtable}
These versions use 44.2 percent fewer simulated instructions in MSA mode. The counts
agree across both byte orders, both tested address layouts and two valid kernel
implementations. Same, all-different, mixed and high-bit numeric patterns give the
same counts for each mode because the baseline forms one Boolean per word without
a data-dependent branch. This is a kernel observation, not a whole-video or elapsed-time
speedup. A vector instruction can also require more simulator host work than a scalar
instruction, so measure wall time separately when that question matters.

A possible next experiment is to bypass four lane extractions when an entire result
vector is zero. Measure both unchanged and all-changing input groups; a new branch
can reduce work on one distribution and add work on another. Retain exact XOR values,
changed counts, alias safety and bounds. Equivalent implementations are not required
to reproduce these particular instruction totals.
